\documentclass[11pt,a4paper]{article}
\usepackage{geometry}
\geometry{margin=1in}
\usepackage{amsmath, amssymb}
\usepackage{graphicx}
\usepackage{booktabs}
\usepackage{hyperref}

\title{\textbf{Hierarchical Model Specifications}\\ \Large Mathematical Formulations for WSLS, CFMR, and ECCM}
\author{Antigravity AI Pipeline}
\date{August 2026}

\begin{document}
\maketitle

\begin{abstract}
This document formally defines the mathematical architectures of the three models evaluated in the Hierarchical Bayesian parameter estimation pipeline. We detail the Win-Stay, Lose-Shift (WSLS) baseline, the Counterfactual Motor Representation (CFMR) baseline, and the ultimate Unified Cortico-Cerebellar Mixed Model (ECCM).
\end{abstract}

\section{The Decision Rule: Drift Diffusion Model (DDM)}
All three models utilize the Drift Diffusion Model (DDM) to generate reaction times (RT) and choices. The probability density function for a choice $C \in \{1, 2\}$ at time $t$ is:
\begin{equation}
    RT \sim \text{Wiener}(v^{(t)}, a, t_{nd}, w=0.5)
\end{equation}
Where $a$ is the decision boundary, $t_{nd}$ is the non-decision time, and $w=0.5$ is the unbiased starting point. The distinguishing feature of each model is how it computes the drift rate $v^{(t)}$ on a trial-by-trial basis.

\section{Model 1: Win-Stay, Lose-Shift (WSLS)}
The WSLS model is a simple probabilistic tracking heuristic. It assumes the subject relies purely on the outcome of the immediately preceding trial. 
\subsection{Intuition}
If you win, how likely are you to stay? If you lose, how likely are you to shift? This requires minimal memory and no gradual value learning.
\subsection{Mathematics}
Let $p_{stay}$ and $p_{shift}$ be the probability parameters. The drift rate towards choice 1 is:
\begin{equation}
v^{(t)} = \beta_v \cdot 2 \times \begin{cases}
    p_{stay} - 0.5 & \text{if previous choice was 1 and rewarded} \\
    0.5 - p_{stay} & \text{if previous choice was 2 and rewarded} \\
    0.5 - p_{shift} & \text{if previous choice was 1 and unrewarded} \\
    p_{shift} - 0.5 & \text{if previous choice was 2 and unrewarded}
\end{cases}
\end{equation}

\section{Model 2: CFMR (Counterfactual Q-Learning)}
The CFMR represents the Cortical probability tracker. It learns the macroscopic reward rates of the environment using an asymmetric learning rate.
\subsection{Intuition}
The brain updates its belief about the chosen option based on whether it won or lost. Because it assumes a zero-sum environment, it counterfactually updates the unchosen option in the exact opposite direction.
\subsection{Mathematics}
Let $Q_{1}$ and $Q_{2}$ be the action values. The drift rate is driven by the value difference:
\begin{equation}
    v^{(t)} = \beta_v \cdot (Q_1^{(t)} - Q_2^{(t)})
\end{equation}
Values are updated using asymmetric Long-Term Potentiation (LTP) and Depression (LTD):
\begin{equation}
Q_{chosen}^{(t+1)} = Q_{chosen}^{(t)} + \begin{cases} 
\eta_{LTP} \cdot (1 - Q_{chosen}^{(t)}) & \text{if rewarded} \\
\eta_{LTD} \cdot (-1 - Q_{chosen}^{(t)}) & \text{if unrewarded}
\end{cases}
\end{equation}

\section{Model 3: ECCM (Parameterized Modulatory Cerebellum)}
The ECCM is the ultimate integration of the CFMR Cortex and a high-dimensional Cerebellar Reservoir ($160 \to 1024 \to 256$).
\subsection{Intuition}
The massive Cerebellar network remembers the exact temporal sequence of the last 80 trials. It independently tries to predict the reward. If the Cerebellar prediction ($Q_{CB}$) agrees with the Cortical belief ($Q_{CTX}$), it multiplicatively amplifies the drift rate (faster reaction). If it disagrees (e.g., detecting a reversal before the Cortex realizes it), it suppresses the cortical signal to prevent a confident mistake.
\subsection{Mathematics}
\textbf{The Deep Orthogonal Memory:} Mossy Fibers ($MF$) hold the last 80 choices and 80 outcomes. Granule Cells ($GC$) project this into a 1024-dimensional space, and Molecular Layer Interneurons ($MLI$) provide 256-dimensional inhibition.
\begin{align}
    gc_j &= \tanh\left( \sum W_{MF \to GC} \cdot MF \right) \\
    mli_k &= \tanh\left( \sum W_{GC \to MLI} \cdot gc \right) \\
    Q_{CB} &= \sum W_{PF} \cdot gc - \sum W_{MLI} \cdot mli
\end{align}
\textbf{Thalamic Multiplicative Modulation:} The Cerebellum projects to the Thalamus to modulate the Cortical drift rate:
\begin{equation}
    v^{(t)} = \beta_v \cdot \left[ Q_{CTX, 1} \cdot (1 + w_{cb} Q_{CB, 1}) - Q_{CTX, 2} \cdot (1 + w_{cb} Q_{CB, 2}) \right]
\end{equation}
This mathematically bounds the prediction and prevents Negative Log-Likelihood blowouts.

\end{document}
